\documentclass[11pt]{article}

\usepackage{amssymb,amsmath,amsfonts,amsthm}
\usepackage{lmodern}
\usepackage{xcolor}
\usepackage[colorlinks=true]{hyperref}
\definecolor{cieblue}{RGB}{31, 119, 180}
\hypersetup{
    citecolor=cieblue,
    allcolors=cieblue
}
\usepackage{mathtools}
\usepackage{enumerate}
\usepackage{booktabs}
\usepackage{float}
\usepackage[ruled]{algorithm2e}
\usepackage{tikz}

\setlength{\textwidth}{15.5cm} \setlength{\headheight}{0.5cm} \setlength{\textheight}{21.5cm}
\setlength{\oddsidemargin}{0.25cm} \setlength{\evensidemargin}{0.25cm} \setlength{\topskip}{0.5cm}
\setlength{\footskip}{1.5cm} \setlength{\headsep}{0cm} \setlength{\topmargin}{0.5cm}

\newtheorem{definition}{Definition}[section]
\newtheorem{proposition}{Proposition}[section]
\newtheorem{lemma}{Lemma}[section]
\newtheorem{theorem}{Theorem}[section]
\newtheorem{corollary}{Corollary}[section]
\newtheorem{remark}{Remark}[section]
\newtheorem{example}{Example}[section]

\newcommand{\Z}{\mathbb{Z}}
\newcommand{\R}{\mathbb{R}}
\newcommand{\N}{\mathbb{N}}
\newcommand{\Pic}{\operatorname{Pic}^0}
\newcommand{\Div}{\operatorname{Div}^0}
\newcommand{\im}{\operatorname{im}}
\newcommand{\tw}{t_{\mathrm{w}}}
\newcommand{\pbar}{\bar{p}}
\newcommand{\norm}[1]{\left\lVert#1\right\rVert}

\allowdisplaybreaks

\begin{document}

\title{Improving the Marked Edge Walk via\\Fiber Decomposition and Gradient-Biased Proposals}
\author{Alaittin K\i rt\i \c{s}o\u{g}lu}
\footnotetext[1]{Department of Applied Mathematics, Illinois Institute of Technology}
\maketitle

\tableofcontents


%═══════════════════════════════════════════════════════════════════════════════
\section{Problems with MEW}\label{sec:problems}
%═══════════════════════════════════════════════════════════════════════════════

We identify two structural inefficiencies in the Marked Edge Walk (MEW)
that worsen with problem size.

\medskip\noindent
\textbf{(P1) Within-fiber waste.}
MEW's cycle basis step picks a non-tree edge $e^+$ uniformly from \emph{all}
non-tree edges. On an $N \times N$ grid with $d = 4$ quadrant districts,
the fraction of non-tree edges that are interior (both endpoints in the same
district) is $f_{\mathrm{int}} = 1 - 1/(N-1)$, which exceeds $80\%$ on a
$6 \times 6$ grid and $96\%$ on $30 \times 30$. Swapping along the fundamental
cycle of an interior edge cannot change the partition. We call these
\emph{within-fiber} steps: they shuffle the spanning tree within the same
equivalence class of compatible trees (Section~\ref{sec:fiber}) without
producing any partition exploration.

For each grid size $N \in \{6, 10, 15, 20, 25, 30\}$, we build the
$N \times N$ grid graph, partition into $d = 4$ quadrants, sample an initial
spanning tree via Wilson's algorithm, and run $10{,}000$ MEW cycle basis steps.
After each step, we record whether the boundary signature (the set of tree
edges crossing district boundaries) changed.

\begin{table}[H]
\centering
\begin{tabular}{crrrrrr}
\toprule
Grid & $n$ & Waste & Speedup & $\log t(G)$ & $\log \prod_i t(G[\xi_i])$ & $\log$ ratio \\
\midrule
$6 \times 6$   &   36 & --\% &  --$\times$ &  -- & -- & -- \\
$10 \times 10$ &  100 & --\% &  --$\times$ &  -- & -- & -- \\
$15 \times 15$ &  225 & --\% &  --$\times$ &  -- & -- & -- \\
$20 \times 20$ &  400 & --\% &  --$\times$ &  -- & -- & -- \\
$25 \times 25$ &  625 & --\% &  --$\times$ &  -- & -- & -- \\
$30 \times 30$ &  900 & --\% &  --$\times$ &  -- & -- & -- \\
\bottomrule
\end{tabular}
\caption{Waste fraction and spanning tree scaling with grid size ($d = 4$ quadrant
districts). Waste is the fraction of MEW cycle basis steps that do not change the
boundary signature; Speedup $= 1/(1-\text{Waste})$. The log ratio
$\log t(G) - \log \prod_i t(G[\xi_i])$ measures the cross-boundary state space
size.}\label{tab:scaling}
\end{table}

\begin{proposition}[Waste fraction scaling]\label{prop:waste-scaling}
For the $N \times N$ grid with $d = 4$ quadrant districts, the waste fraction
$\rho(N)$ satisfies
\[
  \rho(N) \;\to\; 1 \quad \text{as } N \to \infty,
\]
and the empirical data is well-approximated by $\rho(N) \approx f_{\mathrm{int}}(N) = 1 - 1/(N-1)$.
\end{proposition}

\begin{proof}[Heuristic argument]
A cycle basis step picks a uniformly random non-tree edge $e^+$. The probability
that $e^+$ is a boundary edge is at most $|\partial\xi|/|E \setminus T|$, where
$|E \setminus T| = m - (n-1) = g$ is the number of non-tree edges (the genus).
Even if $e^+$ is a boundary edge, the fundamental cycle may swap it for another
boundary edge, leaving the boundary signature unchanged. Conversely, if $e^+$ is
an interior edge, the step can only change the boundary signature if the
fundamental cycle passes through boundary edges and the randomly selected
$e^-$ happens to be one of them.

As $N$ grows, $|\partial\xi| = 2N$ while $g = m - n + 1 = 2N(N-1) - N^2 + 1
= N^2 - 2N + 1 = (N-1)^2$. The fraction of non-tree edges that are boundary
edges is at most $2N/(N-1)^2 \to 0$, so the probability that any given step
touches the boundary vanishes. This explains why the waste fraction approaches
$1$.
\end{proof}

\medskip\noindent
\textbf{(P2) Blind marked-edge selection.}
To understand the relative contribution of the cycle step and the marked edge
step to partition exploration, we classify every MEW step on the $6 \times 6$
grid into four cases based on whether the quotient tree (the set of
boundary tree edges) and the partition change. The quotient tree
captures the cross-boundary structure; two spanning trees with the same quotient
tree differ only in their within-district edges.

We run $50{,}000$ full MEW steps (cycle basis step $+$ marked edge step) on the
$6 \times 6$ grid with $d = 4$ quadrant districts. After each step, we record
the quotient tree and partition.

\begin{table}[H]
\centering
\begin{tabular}{clccr@{\hspace{8pt}}rr}
\toprule
Case & Description & QT & Partition & Count & \% & New discoveries \\
\midrule
A & Boundary jump              & changed & changed & -- & --\% & -- \\
B & Local adjustment (via $M$) & same    & changed & -- & --\% & -- \\
C & QT change, same partition  & changed & same    & -- & --\% & --- \\
D & No change                  & same    & same    & -- & --\% & --- \\
\bottomrule
\end{tabular}
\caption{Classification of MEW steps on the $6 \times 6$ grid by quotient tree (QT)
and partition change. New discoveries counts steps that reach a partition not
previously visited.}\label{tab:step-class}
\end{table}

\begin{enumerate}
  \item Case C = 0: every quotient tree change produces a partition
    change. Boundary moves are never wasted.

  \item Case D $\approx$ 43\%: nearly half of all MEW steps change neither
    the quotient tree nor the partition. These are the within-fiber moves
    of~(P1).

  \item Case A vs.\ B: among steps that change the partition,
    $\sim 85\%$ come from boundary jumps (the cycle step changes the quotient
    tree) and only $\sim 15\%$ from local adjustment (the marked edge step
    slides $M$ within the same quotient tree). The cycle step is the primary
    driver of partition exploration, yet MEW's random walk on marked edges
    treats all tree edges equally.

  \item Discovery rate: boundary jumps discover $\sim 78\%$ of all new
    partitions, but $\sim 80\%$ of boundary jumps revisit previously seen
    partitions. This high revisit rate reflects
    the locality of MEW's single-edge-swap moves.
\end{enumerate}

\begin{remark}[Comparison with ReCom's exploration pattern]
ReCom's merge-and-resplit step completely resamples a spanning tree of the
merged subgraph, producing a partition that can be very different from the
current one. MEW's cycle step swaps one edge at a time, making incremental
changes to the quotient tree.
The advantage of MEW is not mixing speed but the ability to target arbitrary
distributions via calculable transition probabilities~\cite{mew}.
\end{remark}


%═══════════════════════════════════════════════════════════════════════════════
\section{Our Approach}\label{sec:algorithm}
%═══════════════════════════════════════════════════════════════════════════════

We modify the MEW framework to address both problems:
\begin{itemize}
  \item Step~1 (boundary cycle step) solves \textbf{(P1)} by restricting
    $e^+$ to boundary non-tree edges and biasing selection toward edges with
    high population gradient. Every proposal changes the partition.
  \item Step~2 (QP-guided marked edge selection) solves \textbf{(P2)} by
    replacing MEW's blind random walk on marked edges with a quadratic program
    that directly targets population balance.
\end{itemize}

Fix a connected graph $G = (V, E)$ with populations $\{p_v\}$, number of
districts $d$, ideal population $\pbar = \frac{1}{d}\sum_v p_v$, and balance
tolerance $\varepsilon > 0$.

\begin{definition}[Feasible partition]\label{def:feasible}
A $d$-partition $\xi$ is \emph{feasible} if each district is connected and
population-balanced:
\[
  |p(\xi_i) - \pbar| \leq \varepsilon \, \pbar \quad \text{for all } i,
\]
where $p(\xi_i) = \sum_{v \in \xi_i} p_v$.
\end{definition}

The chain operates on the MEW state space $\mathcal{X} = \{(T, M)\}$ where
$T$ is a spanning tree of $G$ and $M \subset E(T)$ with $|M| = d - 1$ is the
marked edge set. The forest $T \setminus M$ defines a $d$-partition
$\xi(T, M)$.

A non-tree edge $e \in E(G) \setminus E(T)$ is a \emph{boundary non-tree edge}
if its endpoints lie in different districts of $\xi$.

\subsection{The Boundary Cycle Step}

The key observation is that MEW's cycle basis step picks $e^+$ uniformly from
\emph{all} non-tree edges. On an $N \times N$ grid, $65\%$--$94\%$ of non-tree
edges are interior (both endpoints in the same district), and swapping along
their fundamental cycle cannot change the partition. Two fixes are applied:
restrict $e^+$ to boundary non-tree edges, and bias selection
toward edges with high population gradient.

\subsubsection*{Pre-computation: Population Potential}

Define the \emph{district imbalance source} for each node $v \in V$ with
$\xi(v) = i$:
\[
  b_v = \frac{p(\xi_i) - \pbar}{|\xi_i|}.
\]
Nodes in overpopulated districts act as sources; nodes in underpopulated
districts act as sinks. The vector $b$ satisfies $\sum_v b_v = 0$.

Solve the Poisson equation $\Delta_G \phi = b$ for the \emph{population
potential} $\phi \in \R^V$. The gradient $|\phi(u) - \phi(v)|$ across each
edge $(u,v)$ measures the inter-district population pressure: edges with large
gradients are natural partition boundaries separating surplus from deficit
districts. For planar redistricting graphs, the sparse Cholesky factorization
of $\Delta_G$ costs $O(n^{3/2})$ and each back-solve is $O(n)$. The potential
$\phi$ is recomputed when the partition changes.

\begin{algorithm}[H]
\caption{MEW with Boundary Cycle Step}\label{alg:boundary-cycle}
\KwIn{State $(T, M)$ with feasible partition $\xi = T \setminus M$;
population potential $\phi$}
\KwOut{New state $(T', M')$}
\BlankLine
\tcc{Step 1: Gradient-biased boundary cycle step}
$B \gets \{e \in E(G) \setminus E(T) : \xi(e_0) \neq \xi(e_1)\}$\;
Select $e^+ \in B$ with $P(e^+) = \dfrac{|\phi(e^+_0) - \phi(e^+_1)|}
  {\sum_{e \in B} |\phi(e_0) - \phi(e_1)|}$\;
$C \gets$ fundamental cycle of $e^+$ in $T$\;
Select $e^- \in C \setminus (M \cup \{e^+\})$ uniformly at random\;
\lIf{no such $e^-$ exists}{\Return{$(T, M)$}}
$T' \gets (T \cup \{e^+\}) \setminus \{e^-\}$\;
\BlankLine
\tcc{Step 2: QP-guided marked edge selection}
\ForEach{$e \in T'$}{
  Compute subtree population $\sigma_e$ and rebalancing flow
  $f_e = \sum_{u \in \mathrm{subtree}(e)} b_u$\;
}
Solve QP: $\min\; x^\top Q\, x + c^\top x$ s.t.\
$\sum x_e = d-1$, $x \in [0,1]^{n-1}$\;
$M' \gets$ $d-1$ edges with largest $x_e$\;
\lIf{$\xi' = T' \setminus M'$ is not feasible}{\Return{$(T, M)$}}
\BlankLine
\tcc{Step 3: Acceptance}
$\alpha \gets \min\!\Bigl(1,\;
  \frac{P'(e^-)}{P(e^+)} \cdot
  \frac{|C \setminus (M \cup \{e^+\})|}{|C' \setminus (M' \cup \{e^-\})|}\Bigr)$\;
Accept $(T', M')$ with probability $\alpha$; otherwise keep $(T, M)$\;
\end{algorithm}

\medskip\noindent
All three components of the MEW transition are modified. Step~1 restricts
$e^+$ to boundary non-tree edges (eliminating within-fiber waste,
Section~\ref{sec:problems}) and biases selection toward high population-gradient
edges via the Laplacian potential $\phi$. Step~2 replaces MEW's blind random
walk on marked edges with QP-guided selection that directly targets population
balance.
Ergodicity over the full state space $\{(T, M)\}$ is preserved: the QP step
can select any subset $M' \subset T'$ of size $d - 1$, and the boundary cycle
step connects all spanning trees reachable via cross-boundary swaps.


\subsection{Key Properties}

\begin{proposition}[Every proposal changes the partition]\label{prop:no-waste}
If $e^+$ is a boundary non-tree edge and $e^- \notin M$, then the partition
$\xi' = T' \setminus M$ is different from $\xi = T \setminus M$.
\end{proposition}

\begin{proof}
Since $e^+$ connects vertices $u \in \xi_i$ and $v \in \xi_j$ with $i \neq j$, the
tree-path from $u$ to $v$ in $T$ must cross at least one marked edge
$m \in M \cap C$.
The edge $e^-$ is a non-marked tree edge, so both its endpoints lie in the same
district $\xi_a$ of $\xi$. Removing $e^-$ from $T$ splits $\xi_a$'s subtree in the
forest $T \setminus M$ into two parts $A_1, A_2$. Adding $e^+$ reconnects one
of these parts (say $A_1$) to a different district via the portion of the cycle
that passes through $e^+$ and possibly other marked edges. Therefore, in
$T' \setminus M$, the vertices of $A_1$ are in a different component than in
$T \setminus M$, and $\xi' \neq \xi$.

Since $M \subset T'$ (we only removed $e^- \notin M$), the forest
$T' \setminus M$ still has exactly $d$ components. Each component is a subtree
of $T'$, hence connected. So $\xi'$ is a valid $d$-partition with connected districts.
\end{proof}

\begin{proposition}[Reversibility]\label{prop:reversible}
The reverse move $e^- \to e^+$ is always a valid boundary cycle step in
state $(T', M)$.
\end{proposition}

\begin{proof}
In state $(T', M)$: $e^-$ is a non-tree edge (we removed it from $T$). Its
endpoints, originally in the same district $\xi_a$ of $\xi$, are now in
\emph{different} districts of $\xi'$ (one part $A_1$ moved to a different
component, as shown above). Therefore $e^-$ is a boundary non-tree edge in
$(T', M)$, so it belongs to $B'$ and can be selected as $e^+$ in the reverse step.

Adding $e^-$ to $T'$ creates a cycle $C'$ that contains $e^+$ (the edge we
previously added). Since $e^+$ was a non-tree edge of $T$ and is now a non-marked
tree edge of $T'$, it belongs to $C' \setminus (M \cup \{e^-\})$ and can be
selected as the removal edge. The reverse swap $(T' \cup \{e^-\}) \setminus \{e^+\} = T$ recovers the original state.
\end{proof}

\begin{proposition}[Correct stationary distribution]\label{prop:stationary}
Algorithm~1 with the acceptance probability $\alpha$ satisfies detailed balance
with respect to the target distribution $\pi(T, M) \propto 1$ on the set of
MEW states with feasible partitions.
\end{proposition}

\begin{proof}[Proof sketch]
The acceptance ratio $\alpha$ is the standard Metropolis--Hastings correction
for the proposal asymmetry. The first factor $|B|/|B'|$ corrects for the
different number of boundary non-tree edges in the two states. The second factor
corrects for the different number of eligible removal edges in the fundamental
cycles. By construction, $\alpha(x \to x') \cdot Q(x \to x') = \alpha(x' \to x)
\cdot Q(x' \to x)$, which is detailed balance.
\end{proof}


\begin{table}[H]
\centering
\begin{tabular}{lcc}
\toprule
Operation & Cost & Notes \\
\midrule
Identify boundary non-tree edges & $O(m)$ & one pass over $E \setminus T$ \\
BFS for fundamental cycle & $O(n)$ & path in $T$ \\
Population balance check & $O(n/d)$ & only changed districts \\
M--H ratio (count $|B'|$, $|C'|$) & $O(m + n)$ & one pass each \\
\bottomrule
\end{tabular}
\caption{Computational cost per step of Algorithm~1.}\label{tab:cost}
\end{table}

\noindent Total cost per step: $O(m + n) = O(n)$ for grid-like graphs --- the
same asymptotic cost as a single MEW cycle basis step. But every step of
Algorithm~1 proposes a genuine partition change, whereas MEW requires
$\sim 1/(1 - \rho)$ steps per useful move, where $\rho$ is the waste fraction
($65\%$ on $6 \times 6$, $93\%$ on $30 \times 30$, $\geq 99\%$ on real graphs).
The effective cost per useful partition move is therefore $O(n) \cdot \frac{1}{1-\rho}$
for MEW versus $O(n)$ for Algorithm~1.


\subsection{QP Details for Marked Edge Selection}\label{sec:qp}

The QP in Step~6 selects the $d - 1$ marked edges by combining two criteria:
the subtree population $\sigma_e$ (cut quality) and the rebalancing flow $f_e$
(partition-aware pressure from the Poisson equation $\Delta_G \phi = b$).

\subsubsection*{Linear coefficients}
For each tree edge $e$:
\[
  c_e = (\sigma_e - \pbar)^2 + (n - \sigma_e - \pbar)^2 - (n - \pbar)^2
  \;-\; \lambda\, f_e^2,
\]
where the first terms penalize imbalanced cuts and $-\lambda f_e^2$ rewards
edges under high rebalancing pressure.

\subsubsection*{Pairwise interactions}
The matrix $Q$ encodes ancestor--descendant dependencies via
inclusion--exclusion on the population imbalance objective
$F(S) = \sum_c (\mathrm{pop}(c) - \pbar)^2$:
\[
  Q_{ef} =
  \begin{cases}
    2\,\sigma_f\,(n - \sigma_e)
      & \text{if } f \text{ is a descendant of } e, \\
    2\,\sigma_e\,(n - \sigma_f)
      & \text{if } e \text{ is a descendant of } f, \\
    2\,\sigma_e\,\sigma_f
      & \text{otherwise (non-ancestor pair).}
  \end{cases}
\]
For ancestor--descendant pairs, $Q_{ef} < 2\sigma_e\sigma_f$ whenever
$\sigma_e > n/2$, discouraging selection of nested edges with large
overlapping subtrees. The pairwise QP is exact for $d \leq 3$ and approximate
for $d \geq 4$ (three-body corrections arise only for ancestor chains).


%═══════════════════════════════════════════════════════════════════════════════
\section{Verification}\label{sec:numerics}
%═══════════════════════════════════════════════════════════════════════════════

We verify the fiber decomposition on a $6 \times 6$ grid with $d = 4$ districts
(quadrants of $3 \times 3$).

\begin{table}[H]
\centering
\begin{tabular}{ll}
\toprule
Quantity & Value \\
\midrule
$|V|, |E|, g$ & $36, 60, 25$ \\
$t(G) = |\Pic(G)|$ & $32{,}565{,}539{,}635{,}200$ \\
$t(G[\xi_i])$ (each quadrant) & $192$ \\
$|\partial(\xi_i, \xi_j)|$ (each adjacent pair) & $3$ \\
$\prod_i t(G[\xi_i])$ & $192^4 = 1{,}358{,}954{,}496$ \\
$\tw(Q_{\xi})$ & $108$ \\
$\tau(\xi) = \prod_i t(G[\xi_i]) \cdot \tw(Q_{\xi})$ & $146{,}767{,}085{,}568$ \\
$\tau(\xi) / t(G)$ & $4.51 \times 10^{-3}$ \\
Boundary configs $|\mathcal{C}_{\xi}|$ & $108$ \\
Trees per boundary config & $1{,}358{,}954{,}496$ \\
\midrule
\multicolumn{2}{l}{\emph{MEW simulation (10{,}000 cycle basis steps):}} \\
Steps changing boundary count & $2{,}799$ $(28.0\%)$ \\
Within-fiber shuffling & $7{,}201$ $(72.0\%)$ \\
\bottomrule
\end{tabular}
\caption{Fiber decomposition verification on the $6 \times 6$ grid.}\label{tab:6x6}
\end{table}


%═══════════════════════════════════════════════════════════════════════════════
\section{Theory}\label{sec:theory}
%═══════════════════════════════════════════════════════════════════════════════

\subsection*{Notation and setup}

Let $G = (V, E)$ be a connected graph with $n = |V|$ vertices and $m = |E|$
edges. The \emph{Laplacian} $\Delta_G$ is the $n \times n$ matrix with
$\Delta_G(v,v) = \deg(v)$ and $\Delta_G(u,v) = -1$ for $u \sim v$. Deleting
the row and column of any vertex $q$ gives the \emph{reduced Laplacian}
$\Delta_G^{(q)}$, and Kirchhoff's matrix-tree theorem states
$t(G) = \det(\Delta_G^{(q)})$, independent of $q$. The \emph{divisor group}
$\operatorname{Div}(G) = \Z^V$ assigns an integer to each vertex; the
degree-zero subgroup is $\Div(G) = \{D \in \Z^V : \sum_v D(v) = 0\}$. The
\emph{Picard group} $\Pic(G) = \Div(G) / \im(\Delta_G)$ satisfies
$|\Pic(G)| = t(G)$, with group structure given by the Smith normal form of
$\Delta_G^{(q)}$. The \emph{genus} is $g = m - n + 1$.


\subsection{Partitions and Compatible Trees}\label{sec:partitions}

\begin{definition}[Partition]\label{def:partition}
A \emph{$d$-partition} of $G$ is a function $\xi : V \to \{1, \ldots, d\}$
such that each \emph{district} $\xi_i = \xi^{-1}(i)$ induces a connected
subgraph $G[\xi_i]$ of $G$. Write $n_i = |\xi_i|$.
\end{definition}

\begin{definition}[Boundary edges]\label{def:boundary}
For a partition $\xi$, the \emph{boundary edge set} between districts $i$ and $j$ is
\[
  \partial(\xi_i, \xi_j) = \bigl\{(u,v) \in E : \xi(u) = i,\; \xi(v) = j\bigr\}.
\]
The full boundary is $\partial\xi = \bigcup_{i < j} \partial(\xi_i, \xi_j)$.
The \emph{interior edge set} of district $i$ is
$E(\xi_i) = \{(u,v) \in E : \xi(u) = \xi(v) = i\}$.
\end{definition}

\begin{definition}[Quotient graph]\label{def:quotient}
The \emph{quotient graph} $Q_{\xi} = (V_Q, E_Q, w)$ is defined as follows:
\begin{itemize}
  \item Vertex set: $V_Q = \{v_1, \ldots, v_d\}$, where $v_i$
    represents district $\xi_i$.
  \item Edge set: $E_Q = \bigl\{(v_i, v_j) : \partial(\xi_i, \xi_j)
    \neq \emptyset\bigr\}$, i.e., two vertices are adjacent if and only if
    their corresponding districts share at least one boundary edge in $G$.
  \item Weights: Each edge $(v_i, v_j) \in E_Q$ carries weight
    $w_{ij} = |\partial(\xi_i, \xi_j)|$, the number of boundary edges between
    $\xi_i$ and $\xi_j$ in $G$.
\end{itemize}
\end{definition}

\begin{definition}[Compatible spanning tree]\label{def:compatible}
A spanning tree $T$ of $G$ is \emph{compatible} with partition $\xi$ if, for
each district $\xi_i$, the restriction $T \cap E(\xi_i)$ is a spanning tree of
$G[\xi_i]$.
\end{definition}

\begin{lemma}[Boundary edge count]\label{lem:boundary-count}
If $T$ is compatible with $\xi$, then $T$ contains exactly $d - 1$ boundary edges,
and these edges form a spanning tree of $Q_{\xi}$ (with each quotient edge realized
by one specific boundary edge).
\end{lemma}

\begin{proof}
$T$ has $n - 1$ edges. The within-district edges account for
$\sum_{i=1}^d (n_i - 1) = n - d$. The remaining $(n-1) - (n-d) = d-1$ edges
are boundary edges. Removing these $d-1$ edges from $T$ disconnects it into
exactly $d$ components (the districts), so the boundary edges must connect all
$d$ districts --- i.e., they form a spanning tree of the quotient graph $Q_{\xi}$.
\end{proof}


\subsection{Fiber Decomposition}\label{sec:fiber}

The central factorization decomposes the set of compatible trees into independent
within-district and cross-boundary choices.

\begin{definition}[Weighted spanning tree count]\label{def:weighted-tree}
For a weighted graph $H$ with edge weights $\{w_e\}$, the \emph{weighted
spanning tree count} is
\[
  \tw(H) = \sum_{T \in \mathcal{T}(H)} \prod_{e \in T} w_e,
\]
where $\mathcal{T}(H)$ is the set of spanning trees of the underlying unweighted
graph. By the weighted matrix-tree theorem, $\tw(H) = \det(\Delta_H^{(q)})$ where
$\Delta_H$ is the weighted Laplacian:
\[
  \Delta_H(i,j) =
  \begin{cases}
    \sum_{j' \neq i} w_{ij'} & \text{if } i = j, \\
    -w_{ij}                   & \text{if } i \sim j, \\
    0                         & \text{otherwise.}
  \end{cases}
\]
\end{definition}

\begin{theorem}[Fiber decomposition]\label{thm:fiber}
The number of spanning trees of $G$ compatible with partition $\xi$ is
\[
  \tau(\xi) = \biggl(\prod_{i=1}^d t\bigl(G[\xi_i]\bigr)\biggr) \cdot \tw(Q_{\xi}),
\]
where $t(G[\xi_i])$ is the unweighted spanning tree count of the $i$-th district
subgraph and $\tw(Q_{\xi})$ is the weighted spanning tree count of the quotient graph.
\end{theorem}

\begin{proof}
A compatible tree $T$ consists of:
\begin{enumerate}[(i)]
  \item A spanning tree $T_i$ of each $G[\xi_i]$, chosen independently
    ($\prod_i t(G[\xi_i])$ choices);
  \item A spanning tree $S$ of $Q_{\xi}$ (specifying which pairs of districts are
    connected), together with a specific boundary edge
    $e_{ij} \in \partial(\xi_i, \xi_j)$ for each edge $(i,j) \in S$
    ($\tw(Q_{\xi})$ choices).
\end{enumerate}
The within-district edges $\bigcup_i E(T_i)$ and the boundary edges
$\{e_{ij} : (i,j) \in S\}$ are disjoint subsets of $E(G)$. Their union has
$(n - d) + (d - 1) = n - 1$ edges, spans all of $V$, and is acyclic (each
$T_i$ is a tree; the quotient tree $S$ is acyclic, so no cycle passes through
multiple districts). Hence the union is a spanning tree of $G$. The map
$(T_1, \ldots, T_d, S, \{e_{ij}\}) \mapsto T$ is a bijection onto the set of
compatible trees.
\end{proof}


\begin{example}[$6 \times 6$ grid, four $3 \times 3$ quadrants]\label{ex:6x6}
\begin{align*}
  t(G) &= 32{,}565{,}539{,}635{,}200, \\
  t\bigl(G[\xi_i]\bigr) &= 192 \quad \text{for each } i = 0,1,2,3, \\
  \prod_i t\bigl(G[\xi_i]\bigr) &= 192^4 = 1{,}358{,}954{,}496.
\end{align*}
The quotient graph $Q_{\xi}$ is a $4$-cycle (districts arranged in a $2 \times 2$
grid). Each pair of adjacent $3 \times 3$ districts shares a boundary along
$3$~nodes, giving $|\partial(\xi_i, \xi_j)| = 3$ boundary edges per adjacent
pair, so all edge weights equal~$3$. The weighted Laplacian is
\[
  \Delta_{Q_{\xi}} =
  \begin{pmatrix}
     6 & -3 & -3 &  0 \\
    -3 &  6 &  0 & -3 \\
    -3 &  0 &  6 & -3 \\
     0 & -3 & -3 &  6
  \end{pmatrix}.
\]
A $4$-cycle has $4$ spanning trees, each using $3$ of $4$ edges. Each tree
contributes $3^3 = 27$, giving $\tw(Q_{\xi}) = 4 \times 27 = 108$.

Verification: $\tw(Q_{\xi}) = \det(\Delta_{Q_{\xi}}^{(0)}) = 108$. \checkmark

The fiber size is
\[
  \tau(\xi) = 192^4 \times 108 = 146{,}767{,}085{,}568,
\]
which is a fraction $\tau(\xi)/t(G) \approx 4.51 \times 10^{-3}$ of all spanning trees.
\end{example}


\begin{definition}[Boundary configuration]\label{def:boundary-config}
For a fixed partition $\xi$, a \emph{boundary configuration} is a pair
$(S, \mathbf{e})$ where $S$ is a spanning tree of $Q_{\xi}$ and
$\mathbf{e} = (e_{ij})_{(i,j) \in S}$ assigns a specific boundary edge
$e_{ij} \in \partial(\xi_i, \xi_j)$ to each edge of $S$.
The set of all boundary configurations is $\mathcal{C}_{\xi}$, with
$|\mathcal{C}_{\xi}| = \tw(Q_{\xi})$.
\end{definition}

\begin{definition}[Within-district equivalence]\label{def:wd-equiv}
Two compatible trees $T, T'$ are \emph{within-district equivalent}
(written $T \sim_{\xi} T'$) if $T \cap \partial\xi = T' \cap \partial\xi$.
Each equivalence class is a \emph{fiber} of size $\prod_i t(G[\xi_i])$.
The number of classes is $|\mathcal{C}_{\xi}| = \tw(Q_{\xi})$.
For the $6 \times 6$ grid: $108$ classes, each containing
$192^4 \approx 1.36 \times 10^9$ trees.
\end{definition}


\subsection{Algebraic Structure: The Quotient Picard Group}\label{sec:picard-quotient}

The fiber decomposition (Theorem~\ref{thm:fiber}) is combinatorial.
We now show that it has a natural algebraic counterpart in the Picard group,
which provides both a conceptual framework and a computational strategy.

The local Picard groups
$\prod_i \Pic(G[\xi_i])$ do not embed into $\Pic(G)$ via chip-firing because
the Laplacian of a subgraph $G[\xi_i]$ is not a sub-matrix of the global
Laplacian (boundary vertex degrees change). However, one can overcome this
obstruction by contracting each district $\xi_i$ into a single
\emph{super-vertex} $v_i$. The resulting graph is exactly the weighted
quotient graph $Q_{\xi}$ from Definition~\ref{def:quotient}, and its Picard group
$\Pic(Q_{\xi})$ captures the cross-boundary structure.

\begin{proposition}\label{prop:quotient-picard}
$|\Pic(Q_{\xi})| = \tw(Q_{\xi})$, and for the $6 \times 6$ grid with four
$3 \times 3$ districts, $|\Pic(Q_{\xi})| = 108$.
\end{proposition}

The group structure is determined by the Smith normal form of
$\Delta_{Q_{\xi}}^{(q)}$. For the $6 \times 6$ grid, the reduced Laplacian is
\[
  \Delta_{Q_{\xi}}^{(0)} =
  \begin{pmatrix}
    6 & 0 & -3 \\
    0 & 6 & -3 \\
    -3 & -3 & 6
  \end{pmatrix},
\]
which has Smith normal form $\operatorname{diag}(3, 36, 0)$, giving
$\Pic(Q_{\xi}) \cong \Z/3 \times \Z/36$ with $|\Pic(Q_{\xi})| = 108$.

Each element of $\Pic(Q_{\xi})$ corresponds to a distinct boundary configuration
$(S, \mathbf{e}) \in \mathcal{C}_{\xi}$ via any standard spanning-tree--to--group
bijection (e.g., Bernardi's bijection or Dhar's burning algorithm).

A degree-zero divisor on $Q_{\xi}$ is a vector $(a_1, \ldots, a_d) \in \Z^d$
with $\sum_i a_i = 0$. To \emph{lift} such a divisor to a divisor on $G$,
choose a representative vertex $u_i \in \xi_i$ for each district and set
\[
  D = \sum_{i=1}^d a_i \cdot [u_i] \;\in\; \Div(G).
\]
This lifting respects linear equivalence: if $D_Q \sim D_Q'$ in $\Pic(Q_{\xi})$
(via chip-firing on $Q_{\xi}$), the lifted divisors $D, D'$ differ by a
chip-firing sequence on $G$ that only fires at the representative vertices.

\begin{remark}[Coset structure]
The lifting defines a map $\Pic(Q_{\xi}) \hookrightarrow \Pic(G)$ whose image
indexes the $\tw(Q_{\xi}) = 108$ cosets of the within-district fiber. Each coset
contains $\prod_i t(G[\xi_i])$ spanning trees that share the same boundary
configuration but differ in their within-district edges.
\end{remark}

To represent each coset uniquely, one can use \emph{$q$-reduced divisors}
(equivalently, $G$-parking functions). Fix a vertex $q \in V$. For each
boundary configuration $(S, \mathbf{e})$, there is a unique $q$-reduced
divisor $D_{(S,\mathbf{e})} \in \Pic(G)$ satisfying:
\begin{enumerate}[(i)]
  \item $D_{(S,\mathbf{e})}(v) \geq 0$ for all $v \neq q$;
  \item no non-empty subset $A \subseteq V \setminus \{q\}$ can be fired
    (i.e., for every such $A$, some $v \in A$ has
    $D_{(S,\mathbf{e})}(v) < \deg_A(v)$, where $\deg_A(v)$ counts edges from
    $v$ to $V \setminus A$).
\end{enumerate}
These $q$-reduced divisors are \emph{stable under within-district
chip-firing}: firing vertices within a single district $\xi_i$ does not
change the $q$-reduced representative, since the within-district moves lie in
the kernel of the projection $\Pic(G) \to \Pic(Q_{\xi})$. Only
\emph{cross-boundary fires} --- those involving vertices in multiple
districts --- change the coset.

\subsection{Computational Implications}

The algebraic structure reduces the complexity of the state space from
$t(G) \approx 3.2 \times 10^{13}$ spanning trees to $|\Pic(Q_{\xi})| = 108$
macro-states. This suggests an algebraic sampling strategy:
\begin{enumerate}
  \item Sample the boundary. Pick a random element from $\Pic(Q_{\xi})$.
    Since $|\Pic(Q_{\xi})| = 108$, this can be done by enumeration or via
    Dhar's burning algorithm on $Q_{\xi}$.
  \item Map to a boundary tree. Use a spanning-tree bijection
    (Bernardi or burning) to convert the group element into a boundary
    configuration $(S, \mathbf{e})$.
  \item Sample the districts. Independently sample a random spanning
    tree $T_i$ of each $G[\xi_i]$ using Wilson's algorithm. Since these are
    small subgraphs (e.g., $3 \times 3$), this is computationally cheap.
  \item Construct the global tree. The union
    $T = (\bigcup_i T_i) \cup \mathbf{e}$ is a compatible spanning tree in the
    selected fiber.
\end{enumerate}
This strategy produces an exact sample from the fiber of any target boundary
configuration, without MCMC. The MCMC component is needed only for exploring
different partitions $\xi$ (via the boundary cycle step of
Section~\ref{sec:algorithm}).

\subsection{The Fiber--Picard Identity}

The algebraic and combinatorial structures are unified by the following
bijection, conditioned on the partition $\xi$:
\[
  \mathcal{T}_{\mathrm{compatible}}(\xi) \;\cong\;
  \biggl(\prod_{i=1}^d \mathcal{T}(G[\xi_i])\biggr) \times \Pic(Q_{\xi}).
\]
Every element of $\Pic(G)$ that projects onto an element of $\Pic(Q_{\xi})$
identifies exactly one fiber. The within-district factor
$\prod_i \mathcal{T}(G[\xi_i])$ lives in the kernel of the projection, and
$\Pic(Q_{\xi})$ indexes the cosets. The fiber decomposition of
Theorem~\ref{thm:fiber} is the counting consequence:
$\tau(\xi) = \prod_i t(G[\xi_i]) \cdot |\Pic(Q_{\xi})|$.


%═══════════════════════════════════════════════════════════════════════════════
\section{Mixing Time}\label{sec:mixing}
%═══════════════════════════════════════════════════════════════════════════════

The boundary cycle step provably improves the spectral gap of the MEW chain.

\begin{proposition}[Peskun dominance]\label{prop:peskun}
Let $P_{\mathrm{MEW}}$ be the transition matrix of the standard MEW chain and
$P_{\mathrm{bnd}}$ that of the boundary cycle chain (Algorithm~1 with Step~1
only, keeping MEW's marked edge walk for Step~2). Both chains share the same
stationary distribution $\pi$ on the state space $\{(T,M)\}$. Then
\[
  P_{\mathrm{bnd}}(x, x) \;\leq\; P_{\mathrm{MEW}}(x, x)
  \quad \text{for all states } x,
\]
and consequently $\gamma_{\mathrm{bnd}} \geq \gamma_{\mathrm{MEW}}$, where
$\gamma$ denotes the spectral gap.
\end{proposition}

\begin{proof}
Both chains use Metropolis--Hastings with the same target $\pi$, so the
stationary distributions agree. The self-loop probability $P(x,x)$ includes
two sources: (i)~the proposal does not change the partition (within-fiber
waste), and (ii)~the proposal is rejected by the M--H step.

In MEW, picking $e^+$ uniformly from all non-tree edges yields within-fiber
waste $\rho \geq 65\%$ (Section~\ref{sec:problems}), so
$P_{\mathrm{MEW}}(x,x) \geq \rho$. In the boundary cycle step, every proposal
changes the partition (Proposition~\ref{prop:no-waste}), so
$P_{\mathrm{bnd}}(x,x)$ equals the rejection probability alone, which is at
most the M--H rejection probability. Since $\rho > 0$,
$P_{\mathrm{bnd}}(x,x) < P_{\mathrm{MEW}}(x,x)$.

By Peskun's theorem~\cite{peskun1973}, if two reversible chains share the same
stationary distribution and $P'(x,x) \leq P(x,x)$ for all $x$, then
$\gamma(P') \geq \gamma(P)$.
\end{proof}


%═══════════════════════════════════════════════════════════════════════════════
\section{Toward an Algebraic Proof of Irreducibility}\label{sec:irreducibility}
%═══════════════════════════════════════════════════════════════════════════════

\subsection{Algebraic Preliminaries}

\begin{definition}[Graph Laplacian]\label{def:laplacian}
Let $G = (V, E)$ be a connected graph with $|V| = n$. The \emph{Laplacian} of
$G$ is the $n \times n$ matrix
\[
  \Delta_G(u,v) =
  \begin{cases}
    \deg(v) & \text{if } u = v, \\
    -1      & \text{if } u \sim v, \\
    0       & \text{otherwise.}
  \end{cases}
\]
The \emph{reduced Laplacian} $\Delta_G^{(q)}$ is obtained by deleting the row
and column of any chosen vertex $q \in V$.
\end{definition}

\begin{definition}[Divisor group]\label{def:divisor}
The \emph{divisor group} of $G$ is $\operatorname{Div}(G) = \Z^V$, the free
abelian group on the vertices. A divisor $D \in \operatorname{Div}(G)$ assigns
an integer $D(v)$ to each vertex. The \emph{degree} of $D$ is
$\deg(D) = \sum_{v \in V} D(v)$. The subgroup of degree-zero divisors is
\[
  \Div(G) = \bigl\{D \in \Z^V : \textstyle\sum_v D(v) = 0\bigr\}.
\]
\end{definition}

\begin{definition}[Chip-firing and linear equivalence]\label{def:chipfire}
\emph{Chip-firing} at vertex $v$ is the map $D \mapsto D - \Delta_G(v, \cdot)$:
vertex $v$ sends one chip along each incident edge to its neighbours. Two
divisors $D, D'$ are \emph{linearly equivalent}, written $D \sim D'$, if
$D - D' \in \im(\Delta_G)$, i.e., $D'$ is reachable from $D$ by a sequence of
chip-fires.
\end{definition}

\begin{definition}[Picard group]\label{def:picard}
The \emph{Picard group} (equivalently, the \emph{sandpile group} or
\emph{critical group}) of $G$ is
\[
  \Pic(G) = \Div(G)\, /\, \im(\Delta_G).
\]
By Kirchhoff's matrix-tree theorem, $|\Pic(G)| = t(G)$, the number of spanning
trees of $G$.
\end{definition}

\subsection{The Graphic Matroid and Its Limitations}

The \emph{graphic matroid} $M(G) = (E, \mathcal{I})$ of $G$ has ground set $E$
(the edges), independent sets $\mathcal{I}$ consisting of all acyclic edge
subsets (forests), and bases equal to the spanning trees of $G$. The
\emph{basis exchange property} states: for any two spanning trees $T, T'$ and
any edge $e \in T \setminus T'$, there exists $f \in T' \setminus T$ such that
$(T \setminus \{e\}) \cup \{f\}$ is again a spanning tree. The MEW cycle basis
step is precisely this operation (swap $e^-$ out, $e^+$ in), and the classical
result that the basis exchange graph is connected gives irreducibility on
spanning trees.

The original MEW extends this to the lifted state space $\{(T, M)\}$ by
combining two facts: (i) any spanning tree is reachable from any other via
basis exchanges, and (ii) any marked set $M'$ is reachable from any $M$ by
sliding marked edges along the tree. Together these imply that all $(T, M)$
pairs are mutually accessible.

However, the graphic matroid sees only the tree $T$; it is blind to the
partition $\xi$ induced by $M$. There is no standard matroid whose bases are
$(T, M)$ pairs, and no matroid-theoretic reason that all partitions should be
reachable. Furthermore, irreducibility on $(T, M)$ does \emph{not} directly
imply irreducibility on the space of feasible (connected, population-balanced)
partitions, because the Metropolis--Hastings rejection step may block
transitions through infeasible intermediates. If the cycle basis step is
restricted to boundary non-tree edges (as in Algorithm~1), the classical
argument (i) itself no longer applies: boundary-restricted swaps cannot
independently rearrange district interiors. Ergodicity of the restricted chain
is therefore an open question.

We propose an algebraic approach to this problem. The Picard group is the
natural algebraic object that sees the partition structure (via cosets of the
projection $\Pic(G) \to \Pic(Q_{\xi})$), and chip-firing---the operation that
generates $\Pic(G)$---may succeed where the matroid argument stops.

\subsection{Chip-Firing as the Elementary Move}

Recall that \emph{chip-firing} at vertex $v$ on a graph $H$ is the map
$D \mapsto D - \Delta_H(v, \cdot)$: vertex $v$ sends one chip along each
incident edge to its neighbours. Two divisors are linearly equivalent if and
only if they are connected by a sequence of chip-fires. The set of chip-fires
at all vertices generates $\Pic(H)$, since the generators of
$\im(\Delta_H)$ are precisely the rows of $\Delta_H$. More precisely, for any
fixed vertex $q$, chip-firing at the $|V|-1$ vertices $v \neq q$ already
suffices, because the row of $q$ is the negative sum of all other rows (the
all-ones vector lies in the kernel of $\Delta_H$).

Applied to our setting, chip-firing operates at two levels:

\begin{enumerate}
  \item \emph{Within-district chip-firing.} Firing a vertex $v$ interior to
    district $\xi_i$ changes the within-fiber representative (the spanning tree
    of $G[\xi_i]$) but preserves the coset in $\Pic(Q_{\xi})$. Algebraically,
    these moves lie in the kernel of the projection
    $\Pic(G) \to \Pic(Q_{\xi})$. They correspond to the within-fiber waste
    identified in Section~\ref{sec:problems}.

  \item \emph{Cross-boundary chip-firing.} Firing a vertex of the quotient
    graph $Q_{\xi}$ (equivalently, simultaneously firing all vertices of a
    district) changes the coset in $\Pic(Q_{\xi})$. Since the chip-fires at
    all $k$ district-vertices generate $\Pic(Q_{\xi})$ (and $k-1$ suffice for
    any fixed sink), all $\tw(Q_{\xi})$ boundary configurations are reachable
    from any starting configuration.
\end{enumerate}

\subsection{Irreducibility for a Fixed Partition}

For a fixed partition $\xi$, the compatible spanning trees are indexed by
$\Pic(G)$ (restricted to the appropriate cosets). Since chip-firing generates
$\Pic(G)$, all compatible trees for a given partition are mutually accessible
via chip-fire sequences. This gives irreducibility within each fiber
\emph{and} across all boundary configurations of a fixed $\xi$.

\subsection{The Open Problem: Changing the Partition}

The difficult step is proving that chip-firing moves can transition between
\emph{different} partitions $\xi \neq \xi'$. Cross-boundary chip-firing on
$Q_{\xi}$ changes the boundary configuration (which edges connect which
districts), but the partition changes only if the new boundary edges alter the
assignment of nodes to districts. Moreover, when $\xi$ changes,
$Q_{\xi}$ itself changes---the algebraic structure shifts with each move.

\begin{enumerate}
  \item[Q1.] Does there exist a chip-firing move on $G$ (or on $Q_{\xi}$) that
    simultaneously (a) has a clean algebraic interpretation as a divisor
    operation and (b) changes the partition $\xi$?

  \item[Q2.] Is the composition of such moves transitive on the space of
    feasible partitions?
\end{enumerate}

An affirmative answer would yield an algebraic proof of irreducibility for the
boundary-restricted chain---something that neither MEW nor ReCom currently
possesses. The Picard group framework reduces the problem to a question about
the action of a finitely generated abelian group on a discrete state space,
bringing tools from algebraic graph theory (Smith normal form, Dhar's burning
algorithm, Bernardi bijections) to bear on a fundamentally combinatorial
question.

\subsection{From Chip-Firing to Tree Moves: the Bijection Gap}\label{sec:bijection-gap}

Chip-firing operates on divisors (integer assignments to vertices), not
directly on spanning trees. The Picard group $\Pic(G)$ has order $t(G)$, so
there is a bijection between its elements and spanning trees---but this
bijection is not canonical. It requires additional structure, and the
combinatorial move that a chip-fire induces on the tree depends entirely on
which bijection is chosen.

We describe three classical bijections and their properties.

\paragraph{Dhar's burning algorithm.}
Fix a sink vertex $q$. A divisor $D$ of degree $g = |E| - |V| + 1$ is
\emph{$q$-reduced} (equivalently, a \emph{$G$-parking function} with respect
to $q$) if:
\begin{enumerate}[(i)]
  \item $D(v) \geq 0$ for all $v \neq q$;
  \item no non-empty subset $A \subseteq V \setminus \{q\}$ can be fired,
    meaning for every such $A$, some $v \in A$ has $D(v) < \deg_A(v)$, where
    $\deg_A(v)$ counts edges from $v$ into $A$.
\end{enumerate}
Every divisor class in $\Pic(G)$ contains a unique $q$-reduced representative.
Dhar's algorithm tests condition (ii) by ``burning'' from $q$: start a fire at
$q$; a vertex $v$ catches fire if the number of its burning neighbours exceeds
$D(v)$; repeat until no new vertex burns. The divisor is $q$-reduced if and
only if all vertices eventually burn. The non-burning vertices identify which
subset to fire to reduce $D$ further.

The bijection with spanning trees is: each $q$-reduced divisor $D$ corresponds
to a unique spanning tree $T_D$, constructed by recording which edges carry the
fire during the burning process. This gives a bijection
$\Pic(G) \leftrightarrow \{\text{spanning trees of } G\}$ that depends on $q$.

\paragraph{Bernardi's bijection.}
Fix a root vertex $q$ and a \emph{ribbon structure} on $G$: a cyclic ordering
of the edges around each vertex (equivalently, an embedding of $G$ into an
orientable surface). Bernardi's bijection constructs a tour of $G$ that
traverses each edge exactly twice (once in each direction), following the
ribbon structure. Given a spanning tree $T$, the tour determines a divisor
$D_T$ by recording the ``break divisor''---the places where the tour must
break away from the tree to follow a non-tree edge.

The key property of Bernardi's bijection is \emph{locality}: if two spanning
trees $T, T'$ differ by a single edge swap (one edge in, one edge out), the
corresponding divisors $D_T, D_{T'}$ differ by a chip-firing move at a
\emph{single} vertex (or a small set of vertices determined by the swap). This
makes Bernardi's bijection the most promising candidate for connecting
chip-firing moves to combinatorial tree operations.

For planar graphs (which includes grid graphs), the ribbon structure is
determined by the planar embedding, removing the choice.

\paragraph{Rotor-routing (cycle-pushing).}
A \emph{rotor configuration} assigns to each non-sink vertex $v$ a distinguished
neighbour $\rho(v)$---the direction of $v$'s ``rotor.'' To route a chip from
$v$: advance $v$'s rotor to the next neighbour in the cyclic order, then move
the chip to that neighbour. Iterate until the chip reaches the sink $q$.

The set of edges $\{(v, \rho(v)) : v \neq q\}$ forms a spanning tree rooted at
$q$ (the rotor tree). Performing one full rotor-routing step changes the rotor
configuration, hence the spanning tree. The rotor-routing process defines a
permutation on spanning trees that corresponds to adding a specific divisor
class in $\Pic(G)$---it is the ``torsor action'' of $\Pic(G)$ on spanning
trees.

\paragraph{The gap.}
All three bijections establish that $|\Pic(G)| = t(G)$ and that chip-firing
acts transitively on divisor classes. However, none of them gives a direct
statement of the form ``chip-firing at boundary vertex $v$ swaps edge $e$ for
edge $f$ in the spanning tree.'' The induced tree move is global: a single
chip-fire changes chips at $v$ and all its neighbours, and the resulting tree
(under any bijection) may differ from the original in multiple edges.

By contrast, the MEW cycle step is a concrete local operation: swap one tree
edge for one non-tree edge. It corresponds to basis exchange in the graphic
matroid---simple and direct, but algebraically opaque.


%═══════════════════════════════════════════════════════════════════════════════
\section{Irreducibility via the Torsor Structure}\label{sec:irreducibility-proof}
%═══════════════════════════════════════════════════════════════════════════════

We now prove that the boundary-restricted chain is irreducible on connected
$k$-partitions. The proof uses the torsor structure of the Picard group acting
on spanning trees, established by Baker and
Wang~\cite{baker2018bernardi} and Backman, Baker, and
Yuen~\cite{backman2019geometric}, combined with the fact that
$\Pic(G)$ is abelian.

\subsection{The Torsor Action}

Let $G$ be a connected planar graph. Baker and Wang~\cite{baker2018bernardi}
prove that the Bernardi process defines a simply transitive action of
$\Pic(G)$ on the set $\mathcal{T}(G)$ of spanning trees:
\[
  \Pic(G) \times \mathcal{T}(G) \to \mathcal{T}(G),
  \qquad (g, T) \mapsto g \cdot T.
\]
Simply transitive means: for any two spanning trees $T, T' \in \mathcal{T}(G)$,
there exists a \emph{unique} group element $g \in \Pic(G)$ such that
$g \cdot T = T'$. The set $\mathcal{T}(G)$ is a \emph{torsor} for $\Pic(G)$.
For planar graphs, this action is canonical---it depends only on the planar
embedding, not on a choice of root vertex
(Theorem~5.1 in~\cite{baker2018bernardi}). Moreover, the Bernardi torsor
coincides with the rotor-routing torsor on planar graphs
(Theorem~7.1 in~\cite{baker2018bernardi}).

\subsection{The Within-District Subgroup}

Let $\xi = (\xi_1, \ldots, \xi_k)$ be a connected $k$-partition of $G$. Define
the \emph{within-district sublattice}
\[
  \Lambda_\xi = \bigl\langle \Delta_G(v, \cdot) : v \text{ is interior to
    some district } \xi_i \bigr\rangle_{\Z} \;\subset\; \im(\Delta_G),
\]
where a vertex $v$ is \emph{interior} to $\xi_i$ if all neighbours of $v$ also
belong to $\xi_i$. Let $\phi : \Z^V \to \Pic(G)$ denote the quotient map.
Define
\[
  H_\xi = \phi(\Lambda_\xi) \;\subset\; \Pic(G).
\]

\begin{lemma}\label{lem:subgroup}
$H_\xi$ is a subgroup of $\Pic(G)$.
\end{lemma}

\begin{proof}
$\Lambda_\xi$ is a sublattice of $\Z^V$ (it is a subgroup of the free abelian
group $\Z^V$ generated by finitely many elements). The map $\phi$ is a group
homomorphism. The image of a subgroup under a group homomorphism is a
subgroup.
\end{proof}

The subgroup $H_\xi$ consists of all elements of $\Pic(G)$ that can be
realised by chip-firing only at interior vertices. Under the torsor action,
$H_\xi$ permutes spanning trees \emph{within the same partition}: if
$T$ is compatible with $\xi$ and $h \in H_\xi$, then $h \cdot T$ is also
compatible with $\xi$, because firing interior vertices rearranges only the
within-district tree edges and does not change the boundary structure.

\begin{lemma}\label{lem:normal}
$H_\xi$ is a normal subgroup of $\Pic(G)$.
\end{lemma}

\begin{proof}
$\Pic(G)$ is a finite abelian group (it is the cokernel of the Laplacian
$\Delta_G$, which is a symmetric matrix). Every subgroup of an abelian group
is normal.
\end{proof}

\subsection{The Quotient Action}

Since $H_\xi$ is normal, the quotient group $\Pic(G) / H_\xi$ is
well-defined. The torsor action of $\Pic(G)$ on $\mathcal{T}(G)$ descends to
an action of $\Pic(G) / H_\xi$ on the set of $H_\xi$-orbits
$\mathcal{T}(G) / H_\xi$:
\[
  \bar{g} \cdot [T]_{H_\xi} = [g \cdot T]_{H_\xi},
  \qquad \bar{g} = gH_\xi \in \Pic(G) / H_\xi.
\]
This is well-defined because $H_\xi$ is normal: if $T' = h \cdot T$ for
$h \in H_\xi$, then $g \cdot T' = g \cdot (h \cdot T) = (gh) \cdot T$, and
$gh$ and $g$ lie in the same coset of $H_\xi$ (since $H_\xi$ is normal, in
fact since $\Pic(G)$ is abelian, $ghH_\xi = gH_\xi$).

\begin{proposition}\label{prop:transitive-orbits}
The action of $\Pic(G) / H_\xi$ on $\mathcal{T}(G) / H_\xi$ is transitive.
\end{proposition}

\begin{proof}
Let $[T]_{H_\xi}$ and $[T']_{H_\xi}$ be two $H_\xi$-orbits. Since
$\Pic(G)$ acts transitively on $\mathcal{T}(G)$ (the torsor property), there
exists $g \in \Pic(G)$ with $g \cdot T = T'$. Then
$\bar{g} \cdot [T]_{H_\xi} = [g \cdot T]_{H_\xi} = [T']_{H_\xi}$.
\end{proof}

\subsection{From Orbits to Partitions}

Each $H_\xi$-orbit consists of spanning trees that share the same
cross-district structure: they are compatible with the same partition and
differ only in the within-district edges. Define the \emph{partition map}
\[
  \pi : \mathcal{T}(G) / H_\xi \;\longrightarrow\; \{\text{connected
    $k$-partitions of } G\}
\]
sending each $H_\xi$-orbit to the partition it represents.

\begin{lemma}\label{lem:surjective}
The map $\pi$ is surjective.
\end{lemma}

\begin{proof}
Every connected $k$-partition $\xi'$ has at least one compatible spanning
tree $T'$ (obtained, for example, by taking a spanning tree of each district
subgraph $G[\xi'_i]$ and connecting them with $k-1$ boundary edges). The orbit
$[T']_{H_\xi}$ maps to $\xi'$ under $\pi$.
\end{proof}

\begin{theorem}[Irreducibility]\label{thm:irreducibility}
Let $G$ be a connected planar graph. The Markov chain on connected
$k$-partitions of $G$ induced by the boundary-restricted cycle step is
irreducible: for any two connected $k$-partitions $\xi$ and $\xi'$, there
exists a sequence of boundary chip-fires (elements of
$\Pic(G) \setminus H_\xi$) transforming a spanning tree compatible with $\xi$
into a spanning tree compatible with $\xi'$.
\end{theorem}

\begin{proof}
Let $\xi$ and $\xi'$ be two connected $k$-partitions. By
Lemma~\ref{lem:surjective}, there exist $H_\xi$-orbits $[T]_{H_\xi}$ and
$[T']_{H_\xi}$ with $\pi([T]_{H_\xi}) = \xi$ and
$\pi([T']_{H_\xi}) = \xi'$. By
Proposition~\ref{prop:transitive-orbits}, there exists
$\bar{g} \in \Pic(G)/H_\xi$ with
$\bar{g} \cdot [T]_{H_\xi} = [T']_{H_\xi}$.

Lift $\bar{g}$ to $g \in \Pic(G)$. Then $g \cdot T$ lies in the orbit
$[T']_{H_\xi}$, which means $g \cdot T$ is compatible with $\xi'$. The
element $g$ is a product of chip-fires (generators of $\Pic(G)$), and since
$\bar{g} \neq \bar{0}$ (assuming $\xi \neq \xi'$), at least one of these
chip-fires lies outside $H_\xi$---i.e., it fires a boundary vertex,
corresponding to a cross-district move.

Since $g$ can be decomposed into a sequence of single chip-fires, each of
which is a valid transition of the boundary-restricted chain (those outside
$H_\xi$) or a within-district move (those inside $H_\xi$, which can be
skipped without affecting the partition), the partition $\xi'$ is reachable
from $\xi$. As $\xi$ and $\xi'$ were arbitrary, the chain is irreducible.
\end{proof}

\begin{remark}
The proof relies on three ingredients:
\begin{enumerate}
  \item The \emph{torsor structure} of $\Pic(G)$ acting on spanning trees,
    established by Baker and Wang~\cite{baker2018bernardi} for planar graphs.
    This guarantees that any spanning tree is reachable from any other by a
    unique group element.
  \item The \emph{abelian} structure of $\Pic(G)$, which ensures that
    $H_\xi$ is normal and the quotient action is well-defined.
  \item The \emph{surjectivity} of the partition map, which ensures every
    partition is represented by at least one orbit.
\end{enumerate}
This gives the first algebraic proof of irreducibility for a redistricting
Markov chain. Unlike the combinatorial proofs used for ReCom~\cite{deford2021recombination} and
flip walks, which rely on ad hoc path-connectivity arguments, the algebraic
proof follows from the group structure of $\Pic(G)$ and applies uniformly to
all connected planar graphs.
\end{remark}



\bibliographystyle{plain}
\bibliography{ref}

\end{document}
